\section{Suficiencia y completitud}
\begin{ejemplo}
    Sean $X_1, X_2$ variables aleatorias independientes tal que $X_1,X_2\rightsquigarrow \cc{P}(\lm)$. Probar que $X_1+X_2$ es suficiente para $\lm$ pero $X_1+2X_2$ no lo es.
    \begin{proof}
        Notamos por $P_\lm$ a la función masa de probabilidad, que dependerá de $\lm$. Queremos probar que, condicionada a un valor de $X_1+X_2$, no depende de $\lm$. Fijado $k\in\bb{N}$, tenemos que:
        \begin{align*}
            P_\lm[X_1=x_1, X_2=x_2 \mid X_1+X_2=k] &= \dfrac{P_\lm[X_1=x_1, X_2=x_2, X_1+X_2=k]}{P_\lm[X_1+X_2=k]}
            =\\&= \begin{cases}
                \dfrac{P_\lm[X_1=x_1, X_2=x_2]}{P_\lm[X_1+X_2=k]} & \text{si } x_1+x_2=k \\
                0 & \text{si } x_1+x_2\neq k
            \end{cases}
        \end{align*}

        Puesto que $X_1$ y $X_2$ son independientes, para $x_1+x_2=k$ tenemos que:
        \begin{align*}
            P_\lm[X_1=x_1, X_2=x_2 \mid X_1+X_2=k] &= \dfrac{P_\lm[X_1=x_1] P_\lm[X_2=x_2]}{P_\lm[X_1+X_2=k]}
            \AstIg \dfrac{e^{-\lm}\frac{\lm^{x_1}}{x_1!} \cdot e^{-\lm}\frac{\lm^{x_2}}{x_2!}}{e^{-2\lm}\frac{(2\lm)^k}{k!}}
            = \\&= \dfrac{k!}{2^k x_1! x_2!}
        \end{align*}
        donde en $(\ast)$ hemos usado la reproductividad de la distribución Poisson, que nos dice que $X_1+X_2\rightsquigarrow \cc{P}(2\lm)$. Notemos que el resultado no depende de $\lm$, por lo que $X_1+X_2$ es suficiente para $\lm$.\\

        Por otro lado, queremos ver que $X_1+2X_2$ no es suficiente para $\lm$. Para ello, no es necesario probarlo para todos los valores posibles, sino que basta con encontrar un contraejemplo. Tomemos $k=2$ y la muestra $(x_1,x_2)=(0,1)$, de forma que $x_1+2x_2=2$. Tenemos que:
        \begin{align*}
            P_\lm[X_1=0, X_2=1 \mid X_1+2X_2=2] &= \dfrac{P_\lm[X_1=0, X_2=1, X_1+2X_2=2]}{P_\lm[X_1+2X_2=2]} =\\&= \dfrac{P_\lm[X_1=0, X_2=1]}{P_\lm[X_1+2X_2=2]}
        \end{align*}

        Por el Teorema de Cambio de Variable, tenemos que:
        \begin{align*}
            P_\lm[X_1+2X_2=2]
            &= \sum\limits_{\substack{x_1,x_2\in\bb{N}\\x_1+2x_2=2}} P_\lm[X_1=x_1, X_2=x_2]
            =\\&= P_\lm[X_1=0, X_2=1] + P_\lm[X_1=2, X_2=0]
        \end{align*}

        Por tanto, tenemos que:
        \begin{align*}
            P_\lm[X_1=0, X_2=1 \mid X_1+2X_2=2] &= \dfrac{P_\lm[X_1=0]P_\lm[X_2=1]}{P_\lm[X_1=0]P_\lm[X_2=1] + P_\lm[X_1=2]P_\lm[X_2=0]} =\\&= \dfrac{e^{-\lm}\frac{\lm^0}{0!} \cdot e^{-\lm}\frac{\lm^1}{1!}}{e^{-\lm}\frac{\lm^0}{0!} \cdot e^{-\lm}\frac{\lm^1}{1!} + e^{-\lm}\frac{\lm^2}{2!} \cdot e^{-\lm}\frac{\lm^0}{0!}} 
            = \dfrac{\lm}{\lm + \frac{\lm^2}{2}} = \dfrac{2}{2+\lm}
        \end{align*}
        Notemos que el resultado depende de $\lm$, por lo que $X_1+2X_2$ no es suficiente para $\lm$.
    \end{proof}
\end{ejemplo}

\begin{ejemplo}
    En $n$ lanzamientos de una moneda, el número de caras es suficiente para la probabilidad de cara (notada $p$).
    \begin{proof}
        Sea $X$ la v.a. que representa el número de caras en un lanzamiento de la moneda, de forma que $X\rightsquigarrow \cc{B}(1,p)$. Sea $(X_1, \ldots, X_n)$ una m.a.s. de $n$ lanzamientos de la moneda, y sea $T = \sum\limits_{i=1}^n X_i$ el número total de caras en los $n$ lanzamientos. Queremos ver que $Y$ es suficiente para $p$. Para ello, fijado un valor $k\in\{0, \ldots, n\}$, tenemos que:
        \begin{align*}
            P_p[X_1=x_1, \ldots, X_n=x_n \mid T=k] &= \dfrac{P_p[X_1=x_1, \ldots, X_n=x_n, T=k]}{P_p[T=k]} =\\&= \begin{cases}
                \dfrac{P_p[X_1=x_1, \ldots, X_n=x_n]}{P_p[T=k]} & \text{si } x_1+\ldots+x_n=k \\
                0 & \text{si } x_1+\ldots+x_n\neq k
            \end{cases}
        \end{align*}

        Por otro lado, como $X_1, \ldots, X_n$ son independientes, en el caso de que se tenga que $x_1+\ldots+x_n=k$, tenemos que:
        \begin{align*}
            P_p[X_1=x_1, \ldots, X_n=x_n \mid T=k] &= \dfrac{P_p[X_1=x_1]\cdots P_p[X_n=x_n]}{P_p[T=k]}
            \AstIg\\&\AstIg \dfrac{\prod\limits_{i=1}^n p^{x_i}(1-p)^{1-x_i}}{\binom{n}{k}p^k(1-p)^{n-k}}
            = \dfrac{p^{\sum\limits_{i=1}^n x_i}(1-p)^{\sum\limits_{i=1}^n 1-x_i}}{\binom{n}{k}p^k(1-p)^{n-k}}
            =\\&= \dfrac{p^k(1-p)^{n-k}}{\binom{n}{k}p^k(1-p)^{n-k}} = \dfrac{1}{\binom{n}{k}}
        \end{align*}
        donde en $(\ast)$ hemos usado que $T\rightsquigarrow \cc{B}(n,p)$. Notemos que el resultado no depende de $p$, por lo que $T$ es suficiente para $p$.
    \end{proof}
\end{ejemplo}

\begin{ejemplo}
    Se supone que se ha lanzado una moneda 100 veces, y se sabe que se han obtenido 60 caras pero se han perdido los datos originales de si salió cara o cruz en cada tirada. Reconstruir la muestra.\\

    Sea $X$ la v.a. que representa el número de caras en un lanzamiento de la moneda, de forma que $X\rightsquigarrow \cc{B}(1,p)$. Sea $(X_1, \ldots, X_{100})$ una m.a.s. de 100 lanzamientos de la moneda, y sea $T = \sum\limits_{i=1}^{100} X_i$ el número total de caras en los 100 lanzamientos, que sabemos que es suficiente para $p$, luego la información que nos da la muestra original se encuentra en $T$. Por tanto, podemos reconstruir la muestra original muestreando de la distribución condicional de $(X_1, \ldots, X_{100})$ dado que $T=60$, que en el ejemplo anterior hemos visto que es:
    \begin{equation*}
        P[X_1=x_1, \ldots, X_{100}=x_{100} \mid T=60] = \begin{cases}
            \dfrac{1}{\binom{100}{60}} & \text{si } x_1+\ldots+x_{100}=60 \\
            0 & \text{si } x_1+\ldots+x_{100}\neq 60
        \end{cases}
    \end{equation*}

    Por tanto, si $(X_1^*, \ldots, X_{100}^*)$ es una m.a.s. de la distribución condicionada al valor de $T=60$, tenemos que $(X_1^*, \ldots, X_{100}^*)$ sigue una distribución uniforme discreta sobre el conjunto de todas las muestras posibles de 100 lanzamientos de la moneda que contienen exactamente 60 caras; es decir, $(X_1^*, \ldots, X_{100}^*)\rightsquigarrow \cc{U}(\cc{S})$, donde:
    \begin{equation*}
        \cc{S} = \left\{(x_1, \ldots, x_{100})\in\{0,1\}^{100} : x_1+\ldots+x_{100}=60\right\}
    \end{equation*}
\end{ejemplo}

\begin{teo}[de Factorización de Neyman-Fisher (Discreto)]
    Sea una v. a. $X\rightsquigarrow F\in \{F_\theta : \theta\in\Theta\}$, y sea $(X_1, \ldots, X_n)$ una m.a.s. de $X$. Sea $f_\theta$ la función masa de probabilidad de $F_\theta$. Sea $T = T(X_1, \ldots, X_n)$ un estadístico. Entonces, $T$ es suficiente para $\theta$ si y solo si existen funciones $g_\theta$ y $h$ tales que:
    \begin{equation*}
        f_\theta(x_1, \ldots, x_n) = h(x_1, \ldots, x_n) g_\theta(T(x_1, \ldots, x_n)) \qquad \forall (x_1, \ldots, x_n)\in\cc{X}^n,\ \forall \theta\in\Theta
    \end{equation*}
    donde $h$ no depende de $\theta$.
    \begin{proof} Demostramos cada una de las implicaciones:
    \begin{description}
        \item[$\Longrightarrow$)] Supongamos que $T$ es suficiente para $\theta$. Entonces, para todo $(x_1, \ldots, x_n)\in\cc{X}^n$, $\theta\in\Theta$, y para cualquier $t\in T(\cc{X}^n)$ fijo, tenemos que:
        \begin{align*}
            P[X_1=x_1, \ldots, X_n=x_n\mid T=t] &= \dfrac{P_\theta[X_1=x_1, \ldots, X_n=x_n, T=t]}{P_\theta[T=t]}
        \end{align*}

        Evaluando en $t=T(x_1, \ldots, x_n)$, tenemos que:
        \begin{align*}
            P[X_1=x_1, \ldots, X_n=x_n\mid T=t] &= \dfrac{P_\theta[X_1=x_1, \ldots, X_n=x_n]}{P_\theta[T=t]}
        \end{align*}

        Despejando, tenemos que:
        \begin{align*}
            P_\theta[X_1=x_1, \ldots, X_n=x_n] &= P_\theta[T=t] P[X_1=x_1, \ldots, X_n=x_n\mid T=t]
        \end{align*}

        Por tanto, identificando
        $$g_\theta(t) = P_\theta[T=t]\qquad \qquad h(x_1, \ldots, x_n) = P[X_1=x_1, \ldots, X_n=x_n\mid T=T(x_1, \ldots, x_n)]$$
        tenemos que:
        \begin{align*}
            P_\theta[X_1=x_1, \ldots, X_n=x_n] = h(x_1, \ldots, x_n) g_\theta(T(x_1, \ldots, x_n))
        \end{align*}

        \item[$\Longleftarrow$)] Supongamos que existen funciones $g_\theta$ y $h$ tales que, para todo $(x_1, \ldots, x_n)\in\cc{X}^n$ y $\theta\in\Theta$, se cumple que:
        \begin{align*}
            P_\theta[X_1=x_1, \ldots, X_n=x_n] = h(x_1, \ldots, x_n) g_\theta(T(x_1, \ldots, x_n))
        \end{align*}

        Entonces, para cualquier $t\in T(\cc{X}^n)$, tenemos que:
        \begin{align*}
            P_\theta[X_1=x_1, \ldots, X_n=x_n\mid T=t] &= \dfrac{P_\theta[X_1=x_1, \ldots, X_n=x_n, T=t]}{P_\theta[T=t]} =\\&= \begin{cases}
                \dfrac{P_\theta[X_1=x_1, \ldots, X_n=x_n]}{P_\theta[T=t]} & \text{si } T(x_1, \ldots, x_n)=t \\
                0 & \text{si } T(x_1, \ldots, x_n)\neq t
            \end{cases}
        \end{align*}

        En el caso de que $T(x_1, \ldots, x_n)=t$, tenemos que:
        \begin{align*}
            P_\theta[X_1=x_1, \ldots, X_n=x_n\mid T=t] &= \dfrac{h(x_1, \ldots, x_n) g_\theta(T(x_1, \ldots, x_n))}{\sum\limits_{\substack{(x_1, \ldots, x_n)\in\cc{X}^n \\ T(x_1, \ldots, x_n)=t}} P_\theta[X_1=x_1, \ldots, X_n=x_n]} =\\&= \dfrac{h(x_1, \ldots, x_n) g_\theta(t)}{\sum\limits_{\substack{(x_1, \ldots, x_n)\in\cc{X}^n \\ T(x_1, \ldots, x_n)=t}} h(x_1, \ldots, x_n) g_\theta(t)} =\\&= \dfrac{h(x_1, \ldots, x_n)}{\sum\limits_{\substack{(x_1, \ldots, x_n)\in\cc{X}^n \\ T(x_1, \ldots, x_n)=t}} h(x_1, \ldots, x_n)}
        \end{align*}
        Notemos que el resultado no depende de $\theta$, por lo que $T$ es suficiente para $\theta$.\qedhere
    \end{description}
    \end{proof}
\end{teo}

\begin{ejemplo}
    Sea $X$ una v.a. tal que $X\rightsquigarrow F\in \{\cc{N}(\mu, \sigma^2) : \mu\in\bb{R}, \sigma^2>0\}$. Estudiar la suficiencia en los casos de:
    \begin{enumerate}
        \item $\sigma_0^2$ conocida.
        
        Consideramos una m.a.s. $(X_1, \ldots, X_n)$ de $X$. Sea $T=T(X_1, \ldots, X_n)$ un estadístico. Por el Teorema de Factorización de Neyman-Fisher, $T$ es suficiente para $\mu$ si y solo si existen funciones $g_\mu$ y $h$ tales que:
        \begin{equation*}
            f_\mu(x_1, \ldots, x_n) = h(x_1, \ldots, x_n) g_\mu(T(x_1, \ldots, x_n)) \qquad \forall (x_1, \ldots, x_n)\in\bb{R}^n,\ \forall \mu\in\bb{R}
        \end{equation*}

        Tenemos que:
        \begin{align*}
            f_\mu(x_1, \ldots, x_n) 
            &\stackrel{\text{i.i.d.}}{=}
            \prod_{i=1}^n f_\mu(x_i)
            = \prod_{i=1}^n \dfrac{1}{\sqrt{2\pi\sigma_0^2}} \exp\left(-\dfrac{(x_i-\mu)^2}{2\sigma_0^2}\right)
            =\\&= \dfrac{1}{(2\pi\sigma_0^2)^{\nicefrac{n}{2}}} \exp\left(-\dfrac{\sum\limits_{i=1}^n (x_i-\mu)^2}{2\sigma_0^2}\right)
            =\\&= \dfrac{1}{(2\pi\sigma_0^2)^{\nicefrac{n}{2}}} \exp\left(-\dfrac{\sum\limits_{i=1}^n (x_i^2 - 2\mu x_i + \mu^2)}{2\sigma_0^2}\right)
            =\\&= \dfrac{1}{(2\pi\sigma_0^2)^{\nicefrac{n}{2}}} \exp\left(-\dfrac{\sum\limits_{i=1}^n x_i^2 - 2\mu \sum\limits_{i=1}^n x_i + n\mu^2}{2\sigma_0^2}\right)
            =\\&= \dfrac{1}{(2\pi\sigma_0^2)^{\nicefrac{n}{2}}} \exp\left(-\dfrac{\sum\limits_{i=1}^n x_i^2}{2\sigma_0^2}\right) \exp\left(-\dfrac{- 2\mu \sum\limits_{i=1}^n x_i + n\mu^2}{2\sigma_0^2}\right)
        \end{align*}

        Por tanto, identificando términos, tenemos que:
        \begin{align*}
            h(x_1, \ldots, x_n) &= \dfrac{1}{(2\pi\sigma_0^2)^{\nicefrac{n}{2}}} \exp\left(-\dfrac{\sum\limits_{i=1}^n x_i^2}{2\sigma_0^2}\right) \\
            g_\mu(T(x_1, \ldots, x_n)) &= \exp\left(-\dfrac{- 2\mu T(x_1, \ldots, x_n) + n\mu^2}{2\sigma_0^2}\right)
        \end{align*}

        Por tanto, por el Teorema de Factorización de Neyman-Fisher, tenemos que $T(X_1, \ldots, X_n) = \sum\limits_{i=1}^n X_i$ es suficiente para $\mu$. En particular, como $\ol{X}= \nicefrac{1}{n} \cdot T$, tenemos que $\ol{X}$ es suficiente para $\mu$.
        \item $\mu_0$ conocida.
        
        En este caso, identificando términos, tenemos que:
        \begin{align*}
            h(x_1, \ldots, x_n) &= 1 \\
            g_{\sigma^2}(T(x_1, \ldots, x_n)) &= \dfrac{1}{(2\pi\sigma^2)^{\nicefrac{n}{2}}} \exp\left(-\dfrac{T(x_1, \ldots, x_n)}{2\sigma^2}\right)
        \end{align*}

        Por tanto, por el Teorema de Factorización de Neyman-Fisher, tenemos que $T(X_1, \ldots, X_n) = \sum\limits_{i=1}^n (X_i-\mu_0)^2$ es suficiente para $\sigma^2$. Como $\mu_0=E[X]=E[\ol{X}]$, tenemos que $S^2 = \nicefrac{1}{n-1} \cdot T$ es suficiente para $\sigma^2$.
        \item $\mu$ y $\sigma^2$ desconocidas.
        
        En este caso, identificando términos, tenemos de forma directa que:
        \begin{align*}
            h(x_1, \ldots, x_n) &= 1
        \end{align*}

        Puesto que el parámetro es bidimensional, el estadístico suficiente debe ser, al menos, bidimensional. Por tanto, podemos tomar como estadístico suficiente $(T_1, T_2)$, donde:
        \begin{equation*}
            T_1(X_1, \ldots, X_n) = \sum\limits_{i=1}^n X_i \qquad \qquad T_2(X_1, \ldots, X_n) = \sum\limits_{i=1}^n (X_i)^2
        \end{equation*}

        De esta forma, tenemos que:
        \begin{multline*}
            g_{\mu, \sigma^2}(T_1(x_1, \ldots, x_n), T_2(x_1, \ldots, x_n)) =\\= \dfrac{1}{(2\pi\sigma^2)^{\nicefrac{n}{2}}} \exp\left(-\dfrac{T_2(x_1, \ldots, x_n) - 2\mu T_1(x_1, \ldots, x_n) + n\mu^2}{2\sigma^2}\right)
        \end{multline*}
        
        De esta forma, por el Teorema de Factorización de Neyman-Fisher, tenemos que $(T_1, T_2)$ es suficiente para $(\mu, \sigma^2)$.
    \end{enumerate}
\end{ejemplo}



\begin{ejemplo}
    Sea $(X_1, \ldots, X_n)$ una m.a.s de la v.a. $X \rightsquigarrow \cc{B}(1, p)$, con $p \in \left]0, 1\right[$. Encontrar un estadístico suficiente y completo asociado a la muestra.\\

    Buscamos un estadístico suficiente. Por el Teorema de Factorización de Neyman-Fisher, $T$ es suficiente para $p$ si y solo si existen funciones $g_p$ y $h$ tales que:
    \begin{equation*}
        P_p[X_1=x_1, \ldots, X_n=x_n] = h(x_1, \ldots, x_n) g_p(T(x_1, \ldots, x_n)) \qquad \forall (x_1, \ldots, x_n)\in\{0, 1\}^n,\ \forall p\in\left]0, 1\right[
    \end{equation*}

    Tenemos que:
    \begin{align*}
        P_p[X_1=x_1, \ldots, X_n=x_n] &\stackrel{\text{i.i.d.}}{=} \prod_{i=1}^n P_p[X_i=x_i] = \prod_{i=1}^n p^{x_i}(1-p)^{1-x_i} = p^{\sum\limits_{i=1}^n x_i}(1-p)^{\sum\limits_{i=1}^n 1-x_i} =\\&= p^{\sum\limits_{i=1}^n x_i}(1-p)^{n-\sum\limits_{i=1}^n x_i}
    \end{align*}

    Por tanto, identificando términos, tenemos que:
    \begin{align*}
        h(x_1, \ldots, x_n) &= 1 \\
        g_p(T(x_1, \ldots, x_n)) &= p^{T(x_1, \ldots, x_n)}(1-p)^{n-T(x_1, \ldots, x_n)}
    \end{align*}

    Por tanto, por el Teorema de Factorización de Neyman-Fisher, tenemos que $T(X_1, \ldots, X_n) = \sum\limits_{i=1}^n X_i$ es suficiente para $p$. Veamos ahora si $T$ es completo. Sea $g$ una función medible tal que $E_p[g(T)] = 0$ para todo $p\in\left]0, 1\right[$. Probemos que $P_p[g(T)=0] = 1$ para todo $p\in\left]0, 1\right[$. Por reproductividad, se tiene $T \rightsquigarrow \cc{B}(n, p)$. Por tanto, para todo $p\in\left]0, 1\right[$, tenemos que:
    \begin{align*}
        E_p[g(T)] &= \sum\limits_{k=0}^n g(k) P_p[T=k] = \sum\limits_{k=0}^n g(k) \binom{n}{k} p^k (1-p)^{n-k}
        =\\&= (1-p)^n \sum\limits_{k=0}^n g(k) \binom{n}{k} \left(\dfrac{p}{1-p}\right)^k = 0
    \end{align*}

    Sea $q = \nicefrac{p}{1-p}$. Entonces, para todo $q\in\left]0, +\infty\right[$, tenemos que:
    \begin{align*}
        \sum\limits_{k=0}^n g(k) \binom{n}{k} q^k = 0
    \end{align*}

    Notemos que el polinomio $\sum\limits_{k=0}^n g(k) \binom{n}{k} q^k$ tiene grado a lo sumo $n$. Por tanto, como tiene infinitas raíces, debe ser el polinomio nulo. Por tanto, tenemos que:
    \begin{align*}
        g(k) \binom{n}{k} = 0 \qquad \forall k\in\{0, \ldots, n\}
    \end{align*}
    Por tanto, tenemos que $g(k) = 0$ para todo $k\in\{0, \ldots, n\}$. Por tanto, hemos visto que:
    \begin{equation*}
        \{0, \ldots, n\} \subset \{k\in\bb{N} : g(k) = 0\}
    \end{equation*}
    Por tanto:
    \begin{equation*}
        1=P_p[T\in \{0, \ldots, n\}] \leq P_p[g(T)=0]\leq 1
    \end{equation*}
    Por tanto, $P_p[g(T)=0] = 1$ para todo $p\in\left]0, 1\right[$; luego $T$ es completo para $p$.
\end{ejemplo}

\begin{ejemplo}
    Sea $(X_1, \ldots, X_n)$ una m.a.s. de la v.a. $X \rightsquigarrow 
    \cc{U}(0, \theta)$, con $\theta > 0$. Encontrar un estadístico suficiente y completo asociado a la muestra.\\
    
    
    Buscamos un estadístico suficiente. Por el Teorema de Factorización de Neyman-Fisher, $T$ es suficiente para $\theta$ si y solo si existen funciones $g_\theta$ y $h$ tales que:
    \begin{equation*}
        f_\theta(x_1, \ldots, x_n) = h(x_1, \ldots, x_n) g_\theta(T(x_1, \ldots, x_n)) \qquad \forall (x_1, \ldots, x_n)\in\bb{R}^n,\ \forall \theta > 0
    \end{equation*}
    Tenemos que:
    \begin{align*}
        f_\theta(x_1, \ldots, x_n) &\stackrel{\text{i.i.d.}}{=} \prod_{i=1}^n f_\theta(x_i) = \prod_{i=1}^n \dfrac{1}{\theta} = \dfrac{1}{\theta^n}\qquad \text{si } (x_1, \ldots, x_n)\in [0, \theta]^n
    \end{align*}

    Notemos que la función de densidad no depende de $(x_1, \ldots, x_n)$, salvo por la condición de que $(x_1, \ldots, x_n)\in [0, \theta]^n$. Por tanto, hemos de emplear la función indicadora para expresar la dependencia de $(x_1, \ldots, x_n)$. Notemos que:
    \begin{align*}
        f_\theta(x_1, \ldots, x_n) & = \dfrac{1}{\theta^n} \ind_{[0,\infty]}(x_{(1)}) \ind_{\left]-\infty, 0\right]}(x_{(n)}-\theta)\qquad \forall (x_1, \ldots, x_n)\in\bb{R}^n,\ \forall \theta > 0
    \end{align*}
    Por tanto, identificando términos, tenemos que:
    \begin{align*}
        h(x_1, \ldots, x_n) &= \ind_{[0,\infty]}(x_{(1)}) \\
        g_\theta(T(x_1, \ldots, x_n)) &= \dfrac{1}{\theta^n} \ind_{\left]-\infty, 0\right]}(T(x_1, \ldots, x_n)-\theta)
    \end{align*} 

    Por tanto, por el Teorema de Factorización de Neyman-Fisher, tenemos que $T(X_1, \ldots, X_n) = X_{(n)}$ es suficiente para $\theta$. Para estudiar su completitud, hemos de calcular en primer lugar la función de densidad de $T$. Tenemos que:
    \begin{align*}
        F_{X_{(n)}}(x) &= P[X_{(n)} \leq x] = P[X_1\leq x, \ldots, X_n\leq x] = \prod_{i=1}^n P[X_i\leq x] = \prod_{i=1}^n F_X(x) = F_X(x)^n
    \end{align*}
    Por tanto, tenemos que:
    \begin{align*}
        f_{X_{(n)}}(x) &= \dfrac{d}{dx} F_{X_{(n)}}(x) = n F_X(x)^{n-1} f_X(x) = n \left(\dfrac{x}{\theta}\right)^{n-1} \dfrac{1}{\theta} = \dfrac{n x^{n-1}}{\theta^n} \qquad \text{si } x\in [0, \theta]
    \end{align*}

    Sea $g$ una función medible tal que $E_\theta[g(T)] = 0$ para todo $\theta > 0$. Probemos que $P_\theta[g(T)=0] = 1$ para todo $\theta > 0$. Tenemos que:
    \begin{align*}
        E_\theta[g(T)] &= \int_0^\theta g(x) f_{X_{(n)}}(x) dx = \int_0^\theta g(x) \dfrac{n x^{n-1}}{\theta^n} dx = \dfrac{n}{\theta^n} \int_0^\theta g(x) x^{n-1} dx = 0\qquad \forall \theta > 0
    \end{align*}

    Por tanto, como el primer término no se puede anular, se tiene:
    \begin{align*}
        \int_0^\theta g(x) x^{n-1} dx = 0\qquad \forall \theta > 0
    \end{align*}
    % // TODO: Por qué es continua??
    Notemos que la función $g(x) x^{n-1}$ es continua, y por tanto integrable. Derivando con respecto a $\theta$, tenemos que:
    \begin{align*}
        g(\theta) \theta^{n-1} = 0\qquad \forall \theta > 0
    \end{align*}
    Esto implica que $g(\theta) = 0$ para todo $\theta > 0$, luego:
    \begin{equation*}
        \bb{R}^+ \subset \{x\in\bb{R} : g(x) = 0\}
    \end{equation*}


    Por tanto:
    \begin{equation*}
        1 = P_\theta[T\in \bb{R}^+ ] \leq P_\theta[g(T)=0]\leq 1
    \end{equation*}
    Por tanto, $P_\theta[g(T)=0] = 1$ para todo $\theta > 0$; luego $T$ es completo para $\theta$.
\end{ejemplo}

\begin{ejemplo}
    Demostrar que las siguientes familias de distribuciones son una familia exponencial uniparamétrica:
    \begin{enumerate}
        \item $\cc{P}(\lm)$, con $\lm > 0$
        
        En primer lugar, tenemos que $\lm\in \Theta = \bb{R}^+$. Además, tenemos que:
        \begin{align*}
            \chi = \bb{N} \qquad \qquad \forall \lm\in\Theta
        \end{align*}

        Por último, para cada $\lm\in\Theta$ y $x\in\chi$, tenemos que:
        \begin{align*}
            P_\lm[X=x] = e^{-\lm}\dfrac{\lm^x}{x!}
            = \exp\left(-\lm + \ln\left(\dfrac{\lm^x}{x!}\right)\right)
            = \exp\left(\ln(\lm)x - \lm - \ln(x!)\right)
        \end{align*}

        Por tanto, identificando términos, tenemos que:
        \begin{align*}
            Q(\lm) &= \ln(\lm) \\
            T(x) &= x \\
            D(\lm) &=- \lm \\
            S(x) &=- \ln(x!)
        \end{align*}

        Por tanto, $\cc{P}(\lm)$ es una familia exponencial uniparamétrica.
        \item $\cc{B}(k_0, p)$, con $p\in\left]0, 1\right[$ y $k_0\in\bb{N}$ fijo.
        
        En primer lugar, tenemos que $p\in \Theta = \left]0, 1\right[$. Además, tenemos que:
        \begin{align*}
            \chi = \{0, 1, \ldots, k_0\}\qquad \qquad \forall p\in\Theta
        \end{align*}

        Por último, tenemos que:
        \begin{align*}
            P_p[X=x] &= \binom{k_0}{x} p^x (1-p)^{k_0-x}
            = \exp\left(\ln\left(\binom{k_0}{x} p^x (1-p)^{k_0-x}\right)\right)
            =\\&= \exp\left(\ln\left(\binom{k_0}{x}\right) + x\ln(p) + (k_0-x)\ln(1-p)\right)
            =\\&= \exp\left(\ln\left(\binom{k_0}{x}\right) + x\ln\left(\dfrac{p}{1-p}\right) + k_0\ln(1-p)\right)
            \qquad \forall x\in\chi,\ \forall p\in\Theta
        \end{align*}

        Por tanto, identificando términos, tenemos que:
        \begin{align*}
            Q(p) &= \ln\left(\dfrac{p}{1-p}\right) \\
            T(x) &= x \\
            D(p) &= k_0\ln(1-p) \\
            S(x) &= \ln\left(\binom{k_0}{x}\right)
        \end{align*}
    \end{enumerate}
    
\end{ejemplo}


\begin{ejemplo}
    La $\cc{N}(\mu, \sigma^2)$ es una familia exponencial uniparamétrica en cada parámetro y bi-paramétrica en los dos parámetros.\\

    Veamos primero que $\cc{N}(\mu, \sigma^2)$ es una familia exponencial uniparamétrica en $\mu$ para $\sigma^2$ conocida. En este caso, tenemos que:
    \begin{align*}
        \Theta &= \bb{R} \\
        \chi &= \bb{R} \\
        f_\mu(x) &= \dfrac{1}{\sqrt{2\pi\sigma^2}} \exp\left(-\dfrac{(x-\mu)^2}{2\sigma^2}\right) = \exp\left(-\dfrac{(x-\mu)^2}{2\sigma^2} - \dfrac{1}{2}\ln(2\pi\sigma^2)\right)
        =\\&= \exp\left(-\dfrac{x^2 - 2\mu x + \mu^2}{2\sigma^2} - \dfrac{1}{2}\ln(2\pi\sigma^2)\right)
        = \exp\left(\dfrac{\mu x}{\sigma^2} - \dfrac{\mu^2}{2\sigma^2} - \dfrac{x^2}{2\sigma^2} - \dfrac{1}{2}\ln(2\pi\sigma^2)\right)
    \end{align*}

    Por tanto, identificando términos, tenemos que:
    \begin{align*}
        Q(\mu) &= \dfrac{\mu}{\sigma^2} \\
        T(x) &= x \\
        D(\mu) &= -\dfrac{\mu^2}{2\sigma^2} \\
        S(x) &= -\dfrac{x^2}{2\sigma^2} - \dfrac{1}{2}\ln(2\pi\sigma^2)
    \end{align*}

    Vemos por tanto que $\cc{N}(\mu, \sigma^2)$ es una familia exponencial uniparamétrica en $\mu$ para $\sigma^2$ conocida. Veamos ahora que $\cc{N}(\mu, \sigma^2)$ es una familia exponencial uniparamétrica en $\sigma^2$ para $\mu$ conocida. En este caso, tenemos que:
    \begin{align*}
        \Theta &= \bb{R}^+ \\
        \chi &= \bb{R} \\
        f_{\sigma^2}(x) &= \exp\left(-\dfrac{(x-\mu)^2}{2\sigma^2} - \dfrac{1}{2}\ln(2\pi\sigma^2)\right)
    \end{align*}

    Por tanto, identificando términos, tenemos que:
    \begin{align*}
        Q(\sigma^2) &= -\dfrac{1}{2\sigma^2} \\
        T(x) &= (x-\mu)^2 \\
        D(\sigma^2) &= -\dfrac{1}{2}\ln(2\pi\sigma^2) \\
        S(x) &= 0
    \end{align*}

    Por tanto, $\cc{N}(\mu, \sigma^2)$ es una familia exponencial uniparamétrica en $\sigma^2$ para $\mu$ conocida. Veamos por último que $\cc{N}(\mu, \sigma^2)$ es una familia exponencial bi-paramétrica en $(\mu, \sigma^2)$. En este caso, tenemos que:
    \begin{align*}
        \Theta &= \bb{R}\times\bb{R}^+ \\
        \chi &= \bb{R} \\
        f_{\mu, \sigma^2}(x) &= \exp\left(\dfrac{\mu x}{\sigma^2} - \dfrac{\mu^2}{2\sigma^2} - \dfrac{x^2}{2\sigma^2} - \dfrac{1}{2}\ln(2\pi\sigma^2)\right)
    \end{align*}

    Por tanto, identificando términos, tenemos que:
    \begin{align*}
        Q_1(\mu, \sigma^2) &= \dfrac{\mu}{\sigma^2} \\
        Q_2(\mu, \sigma^2) &= -\dfrac{1}{2\sigma^2} \\
        T_1(x) &= x \\
        T_2(x) &= x^2 \\
        D(\mu, \sigma^2) &= -\dfrac{\mu^2}{2\sigma^2} - \dfrac{1}{2}\ln(2\pi\sigma^2) \\
        S(x) &= 0
    \end{align*}

    Por tanto, $\cc{N}(\mu, \sigma^2)$ es una familia exponencial bi-paramétrica en $(\mu, \sigma^2)$.
\end{ejemplo}

\begin{teo}
    Si se tiene una v.a. $X \rightsquigarrow \{F_\theta, \theta\in\Theta\}$ con familia de funciones asociadas, $\{f_\theta(x), \theta\in\Theta\}$, siendo f.d.d. o f.m.p. según el caso, donde la familia de distribuciones es exponencial k-paramétrica, entonces la familia de distribuciones asociadas a $(X_1, \ldots, X_n)$ (que es una m.a.s. de $X$) es también exponencial k-paramétrica:
    \begin{equation*}
        f_\theta^n(x_1, \ldots, x_n) = \exp\left[\sum\limits_{h=1}^k Q_h(\theta) \sum\limits_{i=1}^n T_h(x_i) + \sum\limits_{i=1}^n S(x_i) + n D(\theta)\right], \qquad (x_1, \ldots, x_n)\in\chi^n
    \end{equation*}
    y se tiene:
    \begin{enumerate}
        \item El estadístico $\left(\sum\limits_{i=1}^n T_1(X_i), \ldots, \sum\limits_{i=1}^n T_k(X_i)\right)$ es suficiente para $\theta$.
        \item Si $k\leq n$ y el conjunto imagen de la función $Q(\theta) = (Q_1(\theta), \ldots, Q_k(\theta))$ contiene a un abierto de $\bb{R}^k$, el estadístico $\left(\sum\limits_{i=1}^n T_1(X_i), \ldots, \sum\limits_{i=1}^n T_k(X_i)\right)$ es también completo.
    \end{enumerate}
    \begin{proof}
        Comprobemos que $(X_1, \ldots, X_n)$ es una familia exponencial k-paramétrica. Por ser la familia de distribuciones asociadas a $X$ una familia exponencial k-paramétrica, tenemos que $\Theta$ es un intervalo de $\bb{R}^k$. Además, tenemos que el espacio muestral de $(X_1, \ldots, X_n)$ es independiente de $\theta$, y es $\chi^n$. Por último, tenemos que:
        \begin{align*}
            f_\theta^n(x_1, \ldots, x_n) &= \prod_{i=1}^n f_\theta(x_i) = \prod_{i=1}^n \exp\left[\sum\limits_{h=1}^k Q_h(\theta) T_h(x_i) + S(x_i) + D(\theta)\right]
            =\\&= \exp\left[\sum\limits_{i=1}^n \sum\limits_{h=1}^k Q_h(\theta) T_h(x_i) + \sum\limits_{i=1}^n S(x_i) + \sum\limits_{i=1}^n D(\theta)\right]
            =\\&= \exp\left[\sum\limits_{h=1}^k Q_h(\theta) \sum\limits_{i=1}^n T_h(x_i) + \sum\limits_{i=1}^n S(x_i) + n D(\theta)\right]
        \end{align*}

        Por tanto, identificando términos, tenemos que:
        \begin{align*}
            Q_h^n(\theta) &= Q_h(\theta) \\
            T_h^n(x_1, \ldots, x_n) &= \sum\limits_{i=1}^n T_h(x_i) \\
            S^n(x_1, \ldots, x_n) &= \sum\limits_{i=1}^n S(x_i) \\
            D^n(\theta) &= n D(\theta)
        \end{align*}

        Por tanto, $(X_1, \ldots, X_n)$ es una familia exponencial k-paramétrica. Veamos ahora que $\left(\sum\limits_{i=1}^n T_1(X_i), \ldots, \sum\limits_{i=1}^n T_k(X_i)\right)$ es suficiente para $\theta$.
        \begin{align*}
            f_\theta^n(x_1, \ldots, x_n) &= \exp\left[\sum\limits_{h=1}^k Q_h(\theta) \sum\limits_{i=1}^n T_h(x_i) + \sum\limits_{i=1}^n S(x_i) + n D(\theta)\right] =\\&= \underbrace{\exp\left[\sum\limits_{i=1}^n S(x_i)\right]}_{h(x_1, \ldots, x_n)} \underbrace{\exp\left[\sum\limits_{h=1}^k Q_h(\theta) \sum\limits_{i=1}^n T_h(x_i) + n D(\theta)\right]}_{g_\theta\left(\sum\limits_{i=1}^n T_1(x_i), \ldots, \sum\limits_{i=1}^n T_k(x_i)\right)}
        \end{align*}
        
        Por el Teorema de Factorización de Neyman-Fisher, tenemos que el estadístico $\left(\sum\limits_{i=1}^n T_1(X_i), \ldots, \sum\limits_{i=1}^n T_k(X_i)\right)$ es suficiente para $\theta$. Respecto a la complitud, al ser más complicado, no se incluye la demostración.
    \end{proof}
\end{teo}

\begin{ejemplo}
    Sea $X$ una v.a. con distribución $\cc{N}(\mu, \sigma^2)$ con $\mu \in \mathbb{R}$ y $\sigma^2 > 0$. Buscar un estadístico suficiente y completo basado en una muestra de tamaño $n$.\\

    Como hemos visto en el ejemplo anterior, $\cc{N}(\mu, \sigma^2)$ es una familia exponencial bi-paramétrica en $(\mu, \sigma^2)$. Por tanto, por el Teorema anterior, tenemos que la familia de distribuciones asociadas a $(X_1, \ldots, X_n)$ es también una familia exponencial bi-paramétrica. Por tanto, el estadístico:
    \begin{equation*}
        \left(\sum\limits_{i=1}^n X_i, \sum\limits_{i=1}^n X_i^2\right)
    \end{equation*}
    es suficiente para $(\mu, \sigma^2)$. Además, veamos que la imagen de $Q(\mu, \sigma^2)$ contiene un abierto de $\bb{R}^2$. Consideremos la aplicación:
    \Func{Q}{\Theta=\bb{R}\times\bb{R}^+}{\bb{R}^2}{(\mu, \sigma^2)}{\left(\dfrac{\mu}{\sigma^2}, -\dfrac{1}{2\sigma^2}\right)}

    Observemos que al variar $\sigma^2 \in \bb{R}^+$, la segunda componente $Q_2 = \nicefrac{-1}{2\sigma^2}$ recorre todo el intervalo $\bb{R}^-$. Para cualquier valor fijado de $Q_2 \in \bb{R}^-$, al variar $\mu \in \bb{R}$, la primera componente $Q_1 = -2\mu Q_2$ recorre todo $\bb{R}$. Por consiguiente, el conjunto imagen de $Q$ es:
    \begin{equation*}
        Q(\Theta) = \bb{R} \times \bb{R}^-
    \end{equation*}
    que es un abierto de $\bb{R}^2$. Por tanto, por el Teorema anterior, si $n \geq 2$, el estadístico:
    \begin{equation*}
        \left(\sum_{i=1}^n X_i, \sum_{i=1}^n X_i^2\right)
    \end{equation*}
    es también completo para $(\mu, \sigma^2)$.
\end{ejemplo}